% packages.tex — все \usepackage в одном месте
% Порядок важен: языковые пакеты идут первыми после класса.

%% ── Язык и кодировка ─────────────────────────────────────────────────────────
\usepackage{polyglossia}
\setdefaultlanguage{russian}
\setotherlanguage{english}

%% ── Математика ───────────────────────────────────────────────────────────────
\usepackage{amsmath}
\usepackage{amssymb}
\usepackage{newunicodechar}

%% ── Таблицы ──────────────────────────────────────────────────────────────────
\usepackage{longtable}  % таблицы на несколько страниц
\usepackage{tabularx}   % таблицы с авто-шириной колонок

%% ── Диаграммы (TikZ) ─────────────────────────────────────────────────────────
% kaobook already loads xcolor, graphicx, tikz, booktabs, listings, enumitem,
% microtype, todonotes and the caption/hyperref stack.
\usetikzlibrary{arrows.meta, positioning, shapes.geometric, fit, backgrounds}

%% ── Sequence diagrams ────────────────────────────────────────────────────────
\usepackage{pgf-umlsd}

%% ── Цитаты и библиография ────────────────────────────────────────────────────
\usepackage{csquotes}   % правильные кавычки для русского/английского текста

%% ── Библиография ─────────────────────────────────────────────────────────────
\usepackage{kaobiblio}

%% ── Гиперссылки (kaobook подключает hyperref, добавляем настройки) ───────────
\hypersetup{
  pdftitle  = {Платёжный стек},
  pdfauthor = {},
  colorlinks = true,
  linkcolor  = RoyalBlue,
  citecolor  = ForestGreen,
  urlcolor   = Maroon,
}

%% ── Прочее ───────────────────────────────────────────────────────────────────
\usepackage{pdflscape}  % отдельные страницы в альбомной ориентации
